% Custom CV — IFPEN Solaize — Kinetic model discrimination in Fischer-Tropsch synthesis (R12-2) — 2026-11-12
% Base: resumes/library/process_en/source.tex
\documentclass[10pt]{article}
\usepackage[a4]{nermine-cv}
\begin{document}
\cvheader{Nermine Moumen}{Industrial Chemical Engineering Student}{nermine.moumen@insat.ucar.tn}{+21627040414}{linkedin.com/in/nermine-moumen}

\section{EDUCATION}
\entry{National Institute of Applied Sciences and Technology (INSAT)}{\textit{2022–Present}}
\summary{Final-year Industrial Chemical Engineering student, specialized in Process and Environment, with a broad foundation spanning Thermodynamics, Chemical Catalysis, Energy \& Mass Balances and Renewable energy. Seeking a final-year internship in reactor and kinetic modelling for sustainable fuels, starting February 2027.}

\section{PROFESSIONAL EXPERIENCE}
\entry{SEGULA Technologies – Electrochemical Battery Modeling Internship}{\textit{06/2026–08/2026}}
\begin{bul}
\item Modeled a \textbf{Li-ion cell} (NMC811/Graphite-SiOx) using the Single Particle Model (SPM) implemented in MATLAB.
\item Solved the governing PDEs using the \textbf{Finite Volume Method (FVM)} with Crank–Nicolson time integration.
\item Optimized the mesh and fitted the \textbf{diffusion coefficient} values, reducing the RMSE error from \textbf{202.13 mV to 40.5 mV}.
\item Extracted and processed \textbf{experimental discharge data} from a database to validate MATLAB and COMSOL results.
\end{bul}
\entry{Control Union – Quality Control Internship}{\textit{08/2025–09/2025}}
\begin{bul}
\item Analyzed lubricants for viscosity and heavy metals (\textbf{ICP-OES}).
\item Prepared and interpreted \textbf{analytical reports} to verify compliance with ASTM D5185 and \textbf{ISO 17025} standards.
\end{bul}
\entry{ETAP – Water Treatment and Petrochemistry Internship}{\textit{07/2024–08/2024}}
\begin{bul}
\item Studied water treatment methods (\textbf{coagulation-flocculation, ion exchange}) for contaminant and corrosion control.
\item Performed organic geochemistry techniques (GC, GC-MS, Rock-Eval pyrolysis) for hydrocarbon characterization.
\end{bul}
\section{ADDITIONAL EXPERIENCE}
\entry{Cascade – Online Chemistry Teacher (Part-time)}{\textit{2024–Present}}
\begin{bul}
\item Provide individualized online chemistry tutoring for French Terminale students, covering exam-level concepts.
\end{bul}

\section{ACADEMIC PROJECTS}
\project{Ethylene Oxide Production Process Simulation using Aspen HYSYS}
\begin{bul}
\item Modeled a fixed-bed tubular \textbf{gas-solid catalytic reactor} from its operating parameters and geometry, with \textbf{Langmuir-Hinshelwood kinetics} of two reactions and NRTL-RK thermodynamics.
\item Optimized operating conditions via \textbf{sensitivity analysis}, achieving 96\% recovery and \textbf{97.14\% final purity}.
\end{bul}
\project{Techno-Environmental Analysis of Ethylene Production: Steam Cracking vs Bioethanol Dehydration}
\begin{bul}
\item Simulated both routes in \textbf{Aspen Plus V14}, including bioethanol dehydration over an \textbf{HZSM-5 acid catalyst}, comparing reactor design, separation train and energy integration (\textbf{99.95 mol\%} ethylene).
\item Conducted a cradle-to-gate \textbf{Life Cycle Assessment} in OpenLCA using Aspen Plus outputs as inventory data.
\end{bul}
\project{Prediction of Liquid–Liquid Equilibria and Cross-flow and Countercurrent Extraction for Ternary Systems}
\begin{bul}
\item Formulated the LLE thermodynamic system and solved it via Rachford-Rice and \textbf{arc-length continuation}.
\item Reproduced 7 Mathematica notebooks proving equivalence of the Hunter-Nash and McCabe-Thiele methods.
\item Reimplemented workflow in Python, MATLAB, and Mathematica (spline/fsolve, SciPy/mpltern, Interpolation/FindRoot).
\end{bul}
\project{PEM Fuel Cell and Electrolyzer Modeling with COMSOL and ANSYS Fluent}
\begin{bul}
\item Simulated isothermal and nonisothermal \textbf{PEMFC} behavior in COMSOL Multiphysics, achieving a cell voltage of 0.584 V.
\item Modeled a 3D \textbf{PEM electrolyzer} in Ansys Fluent, achieving \textbf{1.40 L/h H$_2$ production} and identifying the oxygen evolution reaction as the dominant kinetic bottleneck.
\end{bul}

\section{TECHNICAL SKILLS}
\begin{bul}
\item \textbf{Simulation:} Aspen Plus, Aspen HYSYS, COMSOL Multiphysics\sep\textbf{Programming:} \textit{MATLAB, Mathematica, Python}
\item \textbf{Data \& Office:} Microsoft Office Suite, OpenLCA\sep\textbf{Standards:} ISO 9001, ISO 14001, ISO 45001, ISO 17025, ISO 14040
\item \textbf{Lab Techniques:} Chromatography (HPLC, GC, TLC), Spectrometry (MS, ICP-OES), Spectroscopy (NMR, FTIR, UV-Vis), Diffraction (XRD), Microscopy (SEM, TEM), Electrochemistry (Cyclic voltammetry, Polarography).
\end{bul}

\section{EXTRACURRICULAR ACTIVITIES}
\begin{bul}
\item \textbf{IEEE INSAT} – Chair of EMBS; Organizing Chair of TRSYP 3.0; Social Activities Chair of MENASYP.
\item \textbf{THEATRO INSAT} – Member of the dance team.
\end{bul}

\section{LANGUAGES AND CERTIFICATES}
\twocol{Languages}{English (Fluent), French (Fluent), Arabic (Native).}{Certificates}{Aspen HYSYS – MATLAB}
\end{document}
